\documentclass[11pt]{article}
\usepackage[margin=1in]{geometry}
\usepackage{amsmath,amssymb,amsthm,mathtools}
\usepackage{hyperref}
\usepackage{enumitem}
\hypersetup{colorlinks=true,linkcolor=blue,urlcolor=blue}
\newcommand{\Z}{\mathbb{Z}}
\newcommand{\R}{\mathbb{R}}
\newcommand{\C}{\mathbb{C}}
\newcommand{\im}{\operatorname{im}}
\newcommand{\rank}{\operatorname{rank}}
\newcommand{\GL}{\operatorname{GL}}
\newcommand{\SL}{\operatorname{SL}}
\newcommand{\Sp}{\operatorname{Sp}}
\newcommand{\Tr}{\operatorname{Tr}}
\DeclareMathOperator{\Span}{span}
\DeclareMathOperator{\Hom}{Hom}
\title{The Cobordism Functor, the Theta Register, and Root Systems\\[4pt]
\large Problems}
\author{}
\date{}
\begin{document}
\maketitle

Problems for \path{docs/design/cobordism_functor_reading.tex}, in the order
of its sections; theorem, proposition and section numbers refer to that
document and are quoted with their names. Problems marked $(\star)$ are the
spine of the framework; the others are context and practice. Hints follow
each section; complete worked solutions are in the companion
\path{cobordism_functor_solutions.tex}. Code exercises use the package
\texttt{tessera} and \path{examples/cobordism/qubit_animation.py}; their
records land under \path{~/cobordism-runs/}.

\part*{The primer}

\section{Cochains, the Laplacian and the harmonic space (Section~2)}

\begin{enumerate}[label=\textbf{P\arabic*.},leftmargin=*,start=1]
\item $(\star)$ Prove the Leibniz rule $d(\alpha\cup\beta) = d\alpha\cup\beta
  + (-1)^{|\alpha|}\alpha\cup d\beta$ for simplicial cochains, and deduce the
  orthogonality half of Theorem~3.1 (mutual annihilators).
\item $(\star)$ The minimal $\Delta$-complex torus has one vertex $v$,
  three loop edges $a,b,c$ and two triangles $U,L$ with $\partial U = a+b-c$,
  $\partial L = -(a+b-c)$. Write $d_0$ and $d_1$, verify $d_1d_0 = 0$, and
  compute the Betti numbers. Now give the faces masses $M_2 = \operatorname{diag}(m_U,m_L)$
  and the edges any invertible diagonal mass $M_1$. Show that when $m_L = -m_U$
  the Hodge Laplacian $L_1$ vanishes identically, so $\dim\ker L_1 = 3\ne b_1$.
  Which of the certificates of Section~26 (closedness on images, coclosedness on
  the chain frame, an independent kernel of the same span) detects the failure,
  and why does Theorem~2.2 (discrete Hodge) not apply?
\item Prove Theorem~2.4 (periods do not move) and explain what the
  verifier's checks J1--J3 measure: under a jitter of the edge lengths the
  monodromy is unchanged, the harmonic representative of a class moves, and
  the difference of the two representatives is exactly a coboundary.
\end{enumerate}

\paragraph{Hints.} P1: check on a single simplex, then extend linearly. P2: $d_0 = 0$ because every edge is a loop; $d_1\alpha = (x,-x)$ with $x = \alpha_a+\alpha_b-\alpha_c$.

\section{Restriction and the monodromy (Sections~3--4)}

\begin{enumerate}[label=\textbf{P\arabic*.},leftmargin=*,start=4]
\item $(\star)$ Write the long exact sequence of $(W,\partial W)$ in degrees
  $0,1,2$ for $W = T^2\times[0,1]$ and, using Lefschetz duality, verify
  $\rank(H^1(W)\to H^1(\partial W)) = 2 = \tfrac12 b_1(\partial W)$.
\item Describe $L_W\subset H^1(M_0)\oplus H^1(M_1)\cong\C^4$ for the
  product collar concretely, and show it is isotropic for $\Omega_\partial = J\oplus(-J)$.
\item $(\star)$ For $2\times2$ matrices, $M^{\mathsf T}JM =
  \det(M)\,J$. Show this fails for $4\times4$ by computing
  $M^{\mathsf T}(J\oplus J)M$ for $M = \sigma\oplus\sigma$, $\sigma$ the swap:
  $\det M = +1$ yet the form is reversed. Conclude that the orientation sign
  of Remark~3.3 must be declared per torus, not derived from $\det M$.
\item Show $M = \varphi^*$ for the collar whose far end is relabelled by a
  simplicial automorphism $\varphi$, and compute $\varphi^*$ for the swap and
  for the Dehn twist $(x,y)\mapsto(x+y,y)$. Which is simplicial on the square
  grid with one diagonal, and why not the other? Then explain how the layered
  flip passes of Section~35.2 realize the twist anyway: what is the flipped grid,
  and why does one tetrahedron per flipped edge suffice?
\item $(\star)$ The code reports, under a $75\%$ complex jitter of every
  squared length, that $M$ moves by $1.6\times10^{-14}$ while the harmonic
  representative $Z P_A^{-1}$ moves by $0.7$ with the difference exactly in
  $\im d_0$. Name the theorem each number measures, and say what a nonzero
  non-exact part would have meant.
\item Theorem~4.2 needs an integral hypothesis. Take $P_A = \operatorname{diag}(1,2)$
  and $P_B = \operatorname{diag}(2,1)$: show $M$ is symplectic but not integral,
  and explain why invertibility of $P_A$ over $\C$ cannot replace the
  lattice hypothesis. Then take the connected sum of a product collar with a
  solid torus $S^1\times D^2$ (along a ball). Show that $H^1(W)$ gains a
  class with zero periods on both boundary tori, that $P_A$ becomes $2\times3$,
  and that $M$ is nevertheless unique (compare the identifiability
  proposition of Section~26).
\end{enumerate}

\paragraph{Hints.} P4: $H^1(W,\partial W)\cong H_2(W)\cong H_2(T^2) = \Z$. P6: $J\sigma = \operatorname{diag}(1,-1)$. P8: Theorem~2.4 and Lemma~2.1.

\section{Choi vectors and composition (Sections~5--6)}

\begin{enumerate}[label=\textbf{P\arabic*.},leftmargin=*,start=10]
\item $(\star)$ Prove that a subspace $L\subset A\oplus B$ is the graph of a
  linear map iff $L\cap(0\oplus B) = 0$ and $L$ projects onto $A$, and that
  the map is then unique. Identify $\dim L$.
\item Write $|M\rangle\!\rangle = \sum_i Me_i\otimes e_i^*$ for $M = I$ and
  $M = \sigma$ in the basis $e_i\otimes e_j^*$ of $\C^2\otimes(\C^2)^*$. Which is the
  maximally entangled Bell vector, and what does that say about the product collar?
\item $(\star)$ Carry out the Mayer--Vietoris argument of
  Theorem~6.1 for two product collars glued end to end, and
  show the composite is a product collar with $M = I\cdot I$.
\item $(\star)$ Prove Theorem~6.1 in general: for $W = W_1\cup_{M_1}W_2$ the
  boundary relation is the composite of relations, $L_W = L_{W_2}\circ L_{W_1}$,
  where $L_2\circ L_1 = \{(x_0,x_2):\exists x_1,\ (x_0,x_1)\in L_1,\ (x_1,x_2)\in L_2\}$.
  Where does exactness of Mayer--Vietoris enter, and why do invisible classes
  (the image of the connecting map from $H^0(M_1)$) not spoil the statement?
\item The verifier's check C1 compares $M(2L,\text{none})$ with
  $M(L,\text{swap})^2$ on \emph{separately seeded} collars. Explain why this
  tests the group law but not a seam. Then explain what the flip passes (F4,
  T9) add --- $T_AT_B$ and $T_BT_A$ stacked on one host read different
  matrices --- and what they still do not test.
\end{enumerate}

\paragraph{Hints.} P10: the projection to $A$ restricted to $L$ is injective iff the first condition holds. P11: $\sum_i e_i\otimes e_i^*$. P13: both inclusions; the connecting map changes the extension, not its restrictions. P14: $\{\pm I,\pm\sigma\}$ is abelian; flip passes live on one complex.

\section{What the metric controls (Section~7)}

\begin{enumerate}[label=\textbf{P\arabic*.},leftmargin=*,start=15]
\item Using Remark~7.1, show that $\Phi^{\mathsf T}G\Phi$ changes
  under a vertex gauge but $\tilde\Phi^{\mathsf T}G\Phi$ does not, and explain
  why the two agree when every phase is zero. (The code measured $0.76$
  against $1.3\times10^{-15}$.)
\item Proposition~7.2 says the pinned-boundary defect
  $D = \|(M^\vee)^{\mathsf T}G_BM-G_A\|_F/\|G_A\|_F$ is invariant under interior
  metric changes. Prove it, say why a relaxation of the bulk therefore cannot
  lower it, and explain why $D = 0$ would still not certify positive quantum
  unitarity of $M$.
\end{enumerate}

\paragraph{Hints.} P15: at zero phases $\tilde\Phi = Z B^{-\mathsf T}$ with $B = \Phi^{\mathsf T}Z$.

\section{Weights, phase space and quantization (Sections~8--9)}

\begin{enumerate}[label=\textbf{P\arabic*.},leftmargin=*,start=17]
\item $(\star)$ For a torus group $T = \mathfrak h/\Lambda$ show that the characters
  are $e^{2\pi i\langle\mu,\cdot\rangle}$ with $\mu\in\Lambda^*$, that the
  weights of $V\otimes V'$ are the sums of weights, and that the weight lattice
  of $T_1\times T_2$ is $\Lambda_1^*\oplus\Lambda_2^*$. Explain why ``all
  characters'' is an infinite-dimensional space and what theta quantization
  adds to obtain a finite register.
\item $(\star)$ Let $\phi$ be a real closed $1$-cochain on a genus-$g$ surface with
  normalized periods $a_i,b_i$ (Section~9). Show that gauge transformations
  do not change $(a,b)$, that large gauge transformations shift them by
  integers, and that $J(\Sigma) = H^1(\Sigma;\R)/H^1(\Sigma;\Z)$ is a real
  torus of dimension $2g$ on which $\sum_i da_i\wedge db_i$ has integral
  periods equal to the intersection form. With $z = b-\Omega a$, show that the
  integer shifts of $(a,b)$ are exactly the lattice $\Z^g+\Omega\Z^g$ acting on
  $z$, so that $z$ is a coordinate on the complex torus $\C^g/(\Z^g+\Omega\Z^g)$.
\item On the torus with one vertex, two edges and one face, the differentials of
  the rank-one local system with holonomies $u,v$ are
  $d_0 = \begin{psmallmatrix}u-1\\v-1\end{psmallmatrix}$ and
  $d_1 = \begin{psmallmatrix}-(v-1)&u-1\end{psmallmatrix}$. Verify $d_1d_0 = 0$
  and show that for $(u,v)\ne(1,1)$ every cohomology group vanishes. Why does
  this show that the state-selection phases must be stored separately from the
  neutral Hodge assembly?
\end{enumerate}

\paragraph{Hints.} P18: $a\mapsto a+e_i$ sends $z\mapsto z-\Omega e_i$.

\section{Theta functions and the tensor register (Sections~10--11)}

\begin{enumerate}[label=\textbf{P\arabic*.},leftmargin=*,start=20]
\item $(\star)$ Prove the periodicity $\theta_j(z+1) = \theta_j(z)$ and the
  quasi-periodicity of Section~10 directly from the sum.
\item $(\star)$ Show that
  $\|\theta_j\|(z) = |\theta_j(z;\tau)|e^{-\pi k(\operatorname{Im}z)^2/\operatorname{Im}\tau}$
  is invariant under both translations. (This is the quantity the record's
  figure draws; $|\theta_j|$ alone is dominated by the growth factor.)
\item Using the argument principle around a fundamental
  parallelogram and the quasi-periodicity factor, show each $\theta_j$ has
  exactly $k$ zeros there. Locate them for $k=2$, $j=0,1$.
\item $(\star)$ Verify the clock and shift formulas of Section~10
  and the commutation $\mathsf Z\mathsf X = e^{2\pi i/k}\mathsf X\mathsf Z$.
\item Prove Proposition~11.1 (factorization), and explain why the same argument
  fails for a non-diagonal $\Omega$.
\item $(\star)$ Derive the normalization of Theorem~11.2 for $g = 1$: with
  $z = x+\tau y$ show that
  $\sqrt{2k}\sqrt{\operatorname{Im}\tau}\int_0^1\!\!\int_0^1\overline{\theta_j}\theta_{j'}\,
  e^{-2\pi k y^2\operatorname{Im}\tau}\,dx\,dy = \delta_{jj'}$.
  Explain why the weight $e^{-2\pi k y^2\operatorname{Im}\tau}$ is the square of
  the factor in P21, and how the argument generalizes to $g$ variables.
\end{enumerate}

\paragraph{Hints.} P21: compute how each factor changes under $z\mapsto z+\tau$. P22: the winding of $\theta_j$ around the boundary of the parallelogram is
the total change of $\arg$; the two vertical sides cancel and the two
horizontal sides differ by the factor $e^{-2\pi i k z}$. P25: integrate in $x$ first; the surviving terms combine into one Gaussian in $n+j/k+y$.

\section{The geometric operator algebra (Section~12)}

\begin{enumerate}[label=\textbf{P\arabic*.},leftmargin=*,start=26]
\item $(\star)$ With $X_r|j\rangle = |j+\epsilon_r\rangle$, $Z_r|j\rangle = \omega^{j_r}|j\rangle$
  and $D_{p,q} = X^pZ^q$, prove $Z^qX^{p'} = \omega^{q\cdot p'}X^{p'}Z^q$ and deduce
  the product and involution laws of the reduced algebra,
  $D_{p,q}D_{p',q'} = \omega^{q\cdot p'}D_{p+p',q+q'}$ and
  $D_{p,q}^\dagger = \omega^{p\cdot q}D_{-p,-q}$. Check associativity from the
  scalar cocycle.
\item $(\star)$ Show $\Tr(D_{p,q}^\dagger D_{p',q'}) = d\,\delta_{pp'}\delta_{qq'}$ with $d = k^g$,
  conclude that the $d^2$ operators $D_{p,q}$ form an orthogonal basis of
  $\operatorname{End}(\mathcal Q_k)$ with inverse formula
  $T = \frac1d\sum_{p,q}\Tr(D_{p,q}^\dagger T)D_{p,q}$, and derive the twisted
  convolution $(a\star b)_{r,s} = \sum_{p,q}a_{p,q}b_{r-p,s-q}\,\omega^{q\cdot(r-p)}$.
  Verify it by hand for $k = 2$, $g = 1$ on $a = e_{1,0}$ ($X$), $b = e_{0,1}$ ($Z$):
  compute $a\star b$ and $b\star a$ and compare with $XZ$ and $ZX$.
\item Corollary~12.3 identifies states with positive symbols of geometric trace
  one, $\operatorname{tr}_G(a) = d\,a_{0,0}$. For $k = 2$, $g = 1$ find the
  symbol of $|0\rangle\langle0|$ and of the effect $e = |+\rangle\langle+|$,
  check $\operatorname{tr}_G$ of the state is $1$, and compute the Born
  probability $\operatorname{tr}_G(ae)$ without forming any matrix.
\item Prove Theorem~12.4 in coordinates: $D_{(p_1,p_2),(q_1,q_2)} = D_{p_1,q_1}\otimes D_{p_2,q_2}$
  in the factorized basis, the vectors $|D_{p,q}\rangle\!\rangle/\sqrt d$ are
  orthonormal, and contracting the dual intermediate factor of
  $|T_2\rangle\!\rangle\otimes|T_1\rangle\!\rangle$ gives $|T_2T_1\rangle\!\rangle$.
  Do the last step with matrix units $E_{ij}$ and explain why intermediate
  renormalization would lose information.
\end{enumerate}

\paragraph{Hints.} P26: $Z_rX_s = \omega^{\delta_{rs}}X_sZ_r$ on basis vectors. P27: a nonzero shift has zero diagonal; a nonconstant character sums to zero. P28: $|0\rangle\langle0| = (I+Z)/2$, $|+\rangle\langle+| = (I+X)/2$.

\section{The modular action and the Weil representation (Section~13)}

\begin{enumerate}[label=\textbf{P\arabic*.},leftmargin=*,start=30]
\item $(\star)$ Derive $\rho_k(T) = \operatorname{diag}(e^{\pi i j^2/k})$ for
  $k$ even, and show what goes wrong for $k$ odd.
\item $(\star)$ Derive $\rho_k(S)$ by Poisson summation: apply
  $\sum_n f(n) = \sum_m \hat f(m)$ to the Gaussian in the definition of
  $\theta_j(z/\tau;-1/\tau)$. Recover the Hadamard gate at $k=2$. Show that the
  identity is exact with $j(S,\tau) = (-i\tau)^{1/2}$ and no further constant,
  and that at $k = 4$ the sign in $e^{-2\pi ijl/k}$ is forced (the conjugate
  matrix fails).
\item Show $\rho_k(S)^2$ is the parity operator $j\mapsto -j$ (up to a
  phase), and $(\rho_k(S)\rho_k(T))^3\propto\rho_k(S)^2$, matching
  $S^2 = -I$ and $(ST)^3 = S^2$ in $\SL(2,\Z)$.
\item $(\star)$ Verify the shear formula: for $\Omega\mapsto\Omega+B$,
  $\theta_j$ picks up $e^{\pi i j^{\mathsf T}Bj/k}$ when $k$ is even and $B$
  symmetric integer. Compute it for $B = \begin{psmallmatrix}0&1\\1&0\end{psmallmatrix}$,
  $k = 2$, and identify controlled-$Z$.
\item $(\star)$ Prove Proposition~13.3 (products), and compute the operator
  Schmidt rank of $\rho_2(\text{shear})$ across the two qubits.
\item Explain why $\rho_k(\mathsf M)$ cannot depend on $\tau$, given that the
  transformation law holds for all $\tau$ and the functions $\theta_l(\cdot;\tau)$
  are linearly independent. (The code measures the fitted matrices at two
  moduli to agree to $9\times10^{-16}$, check T7.)
\item $(\star)$ Prove Proposition~13.1 (the coordinate bridge). A period map
  $M_{\rm per} = \begin{psmallmatrix}P&Q\\R&S\end{psmallmatrix}$ acts on $(a,b)$.
  Show that the new normalized period matrix is $\Omega' = (R+S\Omega)(P+Q\Omega)^{-1}$,
  that $z' = b'-\Omega'a' = (P+Q\Omega)^{-\mathsf T}z$, and that the standard
  Siegel matrix acting by $\Omega' = (A\Omega+B)(C\Omega+D)^{-1}$ is
  $\mathsf M = \begin{psmallmatrix}S&R\\Q&P\end{psmallmatrix}$. Specialize to genus
  one and compare with the measured far modulus of the flipped torus (F3),
  $\tau' = (M_{21}+M_{22}\tau)/(M_{11}+M_{12}\tau)$.
\item Prove Proposition~13.4 (compatible geometric transport is unitary): the
  clock operators have a common eigenvector whose shift orbit spans an
  irreducible representation of the finite Heisenberg group; a compatible
  relabelling gives a second such representation, and the intertwiner is
  unique up to a scalar, which adjoint-preservation makes a phase.
\end{enumerate}

\paragraph{Hints.} P30: $e^{\pi i k n^2} = (-1)^{kn}$. P31: the Fourier transform of $e^{\pi i k(n+j/k)^2\tau}$ in $n$ is again Gaussian with
$\tau\mapsto-1/\tau$; the characteristic $j$ becomes a phase $e^{-2\pi i jm/k}$. P36: $\Omega'(P+Q\Omega) = R+S\Omega$ and symplecticity give $S-\Omega'Q = (P+Q\Omega)^{-\mathsf T}$. P37: Schur's lemma for the algebra spanned by the $D_{p,q}$, which is all of $\operatorname{End}(\mathcal Q_k)$ (P27).

\section{Orientation reversal and coherent states (Sections~14--15)}

\begin{enumerate}[label=\textbf{P\arabic*.},leftmargin=*,start=38]
\item Show $\overline{\theta_j(z;\Omega)} = \theta_j(-\bar z;-\bar\Omega)$ and
  hence that orientation reversal acts antiunitarily. For the torus period
  swap $M_{\rm per} = \sigma$, form $N = M_{\rm per}R_{\rm per}$ with
  $R_{\rm per} = \operatorname{diag}(1,-1)$, check that $N$ is symplectic,
  identify $N$ and its Siegel matrix $\mathsf N$, and write the level-two action
  as a matrix composed with conjugation.
\item $(\star)$ Prove Proposition~15.2: if $F_{\Omega'}(z') = f\,\rho\,F_\Omega(z)$ with
  $\rho$ unitary and $f\ne0$, then the normalized conjugate evaluation vectors
  satisfy $c_{\Omega'}(z') = (\bar f/|f|)\,\bar\rho\,c_\Omega(z)$. Explain why
  fitting $\rho$ on value columns and applying it unconjugated to coherent
  states is wrong, and why the two conventions are equally complete.
\item $(\star)$ Prove Theorem~15.3: for an elliptic curve $E$ and a degree-two line
  bundle $L^2$, the two basis sections have no common zero, the coherent-state
  map $E\to\mathbb{CP}^1$ is onto and of degree two, and conjugation preserves
  surjectivity. Then show that at fixed polarization a single elliptic curve
  cannot cover all qutrit rays.
\item Show that for $\Omega = \operatorname{diag}(\tau_1,\tau_2)$ the coherent
  state factorizes, $c_\Omega(z) = c_{\tau_1}(z_1)\otimes c_{\tau_2}(z_2)$, so its
  entanglement entropy in the factorized basis is zero, and explain why a
  nonzero $\Omega_{12}$ can make it nonzero even though the monodromy of the
  tube host is the identity (G2--G4).
\end{enumerate}

\paragraph{Hints.} P40: Riemann--Roch on genus one: $h^0(D) = \deg D$ for $\deg D>0$.

\section{Seams, factorization and the Clifford limit (Sections~16--17)}

\begin{enumerate}[label=\textbf{P\arabic*.},leftmargin=*,start=42]
\item $(\star)$ Write out Mayer--Vietoris for the sphere join and prove
  Proposition~16.1. The code measured the off-block norm of the
  four-torus monodromy at $5\times10^{-15}$ (check D1); which map in the
  sequence is that number testing?
\item $(\star)$ Do the same for a tube identifying a cycle $a_1$ on one torus
  with a cycle $a_2$ on the other, with overlap retracting to $S^1$. Show the restriction subspace satisfies
  $p_1(a_1) = p_2(a_2)$, write the simplest monodromy consistent with it,
  and compute its $\rho_2$ and Schmidt rank.
\item Why does attaching a $1$-handle (a solid tube glued along two disks)
  \emph{not} create a shared class, while gluing along an annulus does?
  Reconcile this with the tube host of Section~35.2: its monodromy is the
  identity (G2, ``the tube alone couples no cycle'') while the far surface's
  period matrix has $\Omega_{12}\ne0$ (G3). Which of the two is topological
  and which is metric?
\item $(\star)$ Compute $CZ|{+}{+}\rangle$ and its reduced state on the first qubit.
  Show that $\mathrm{SWAP}$ maps every product state to a product state
  although its operator Schmidt rank is four. Define the product-input
  entangling witness used in E5 and explain why operator Schmidt rank alone
  (Proposition~13.3) is not it.
\item $(\star)$ Show that the Weil representation normalizes the
  Heisenberg--Weyl group: $\rho(S)\mathsf X\rho(S)^{-1}\propto\mathsf Z$. Why
  does this make its image the Clifford group, and why does Gottesman--Knill
  then say the theta register alone is classically simulable?
\item Reduction modulo $2$ maps $\Sp(4,\Z)\to\Sp(4,\mathbb F_2)\cong S_6$, and the
  two-qubit local Cliffords with the swap map onto a subgroup isomorphic to
  $S_3\wr\Z/2$. Use orders to show that no gate whose symplectic image has
  order five modulo $2$ is a local Clifford or a swap times one, and verify
  the hypotheses for the decagon monodromy
  $M = \begin{psmallmatrix}-1&1&0&0\\-1&0&1&1\\0&0&0&1\\-1&0&0&0\end{psmallmatrix}$:
  $\det M = 1$, $M^{\mathsf T}JM = J$, $M^5 = I$, and $M\bmod2$ has order five.
\end{enumerate}

\paragraph{Hints.} P43: the kernel of $H^1(W_1)\oplus H^1(W_2)\to H^1(S^1)$. P44: $W_1\cup\text{handle}\cup W_2\simeq W_1\vee W_2$. P46: $\mathsf X$ and $\mathsf Z$ are translations by $\tau/k$ and $1/k$;
$S$ exchanges the roles of $1$ and $\tau$. P47: $|S_3\wr\Z/2| = 72$.

\section{From a Lie algebra to its root system; the Weyl group (Sections~18--19)}

\begin{enumerate}[label=\textbf{P\arabic*.},leftmargin=*,start=48]
\item $(\star)$ From $[h,e] = 2e$, $[h,f] = -2f$, $[e,f] = h$
  alone, show that a finite-dimensional irreducible $\mathfrak{sl}(2)$-module
  with highest $h$-eigenvalue $\lambda$ has eigenvalues
  $\lambda,\lambda-2,\dots,-\lambda$, each once, so that $\lambda\in\Z_{\ge0}$
  and the dimension is $\lambda+1$. Identify $\alpha$, $\alpha^\vee$,
  $\omega$, $P$, $Q$, $P/Q$ and $W$ in this notation, and check
  $|\omega|^2 = \tfrac12$ when $|\alpha|^2 = 2$.
\item $(\star)$ For $\mathfrak{sl}(3,\C)$ with $\mathfrak h$
  the traceless diagonal matrices, show that the roots are
  $\varepsilon_i-\varepsilon_j$ with root vectors $E_{ij}$; take
  $\alpha_1 = \varepsilon_1-\varepsilon_2$, $\alpha_2 = \varepsilon_2-\varepsilon_3$
  and compute the Cartan matrix, the angle between the simple roots, the
  highest root, $\omega_1,\omega_2$ in terms of $\alpha_1,\alpha_2$, $P/Q$,
  $\rho$, $h^\vee$, and $W$ as a group of permutations. Draw the six roots
  and the weight lattice and mark $P^+$.
\item Write the Cartan matrix and the Gram matrix of the
  simple roots for $A_1\times A_1$, $A_2$, $B_2$, $G_2$ with long roots of
  length $\sqrt2$. Show that the root lattices of $A_1\times A_1$ and $B_2$
  are both square and those of $A_2$ and $G_2$ both hexagonal (up to scale
  and rotation), and that $W(A_1\times A_1)\subset W(B_2)$ and
  $W(A_2)\subset W(G_2)$ have index two, the missing generator being the
  reflection in a short root.
\item $(\star)$ Prove Proposition~18.3: a Euclidean lattice generated by a triangle with
  side lengths $\sqrt2,\sqrt2,\sqrt6$ has Gram matrix
  $\begin{psmallmatrix}2&-1\\-1&2\end{psmallmatrix}$, its six shortest vectors
  form $A_2$, their reflections preserve the lattice and generate $S_3$, and
  the dual lattice is $P = A^{-1}\Z^2$ with $|P/Q| = 3$. Which of these facts
  uses only scalar length data, and what would a boundary triangulation have
  to satisfy for its harmonic mass Gram to equal $A$?
\item $(\star)$ From the Weyl character formula (Theorem~19.1) derive
  $\operatorname{ch}V_\lambda = x^{\lambda}+x^{\lambda-2}+\dots+x^{-\lambda}$
  for $\mathfrak{su}(2)$, and from the dimension formula derive
  $\dim V_{(a,b)} = \tfrac12(a+1)(b+1)(a+b+2)$ for $\mathfrak{su}(3)$; check
  $(1,0)\mapsto3$, $(1,1)\mapsto8$, $(3,0)\mapsto10$, $(2,2)\mapsto27$.
\item $(\star)$ Prove Proposition~19.2 (Brauer--Klimyk) from the
  character formula. Apply it to $V_1\otimes V_\lambda$ and to
  $V_2\otimes V_0$ for $\mathfrak{su}(2)$, and to
  $\mathbf3\otimes\mathbf3$ and $\mathbf3\otimes\bar{\mathbf3}$ for
  $\mathfrak{su}(3)$, exhibiting a term a wall kills and a term a reflection
  folds with a sign.
\item Verify the Weyl denominator identity
  $\sum_w\varepsilon(w)e^{w\rho} = e^{\rho}\prod_{\alpha>0}(1-e^{-\alpha})$
  by direct expansion for $A_1$ and for $A_2$ (six terms against eight, two
  of which cancel), identifying each surviving term with an element of $W$
  and its sign.
\end{enumerate}

\paragraph{Hints.} P48: $e$ raises and $f$ lowers the
$h$-eigenvalue by $2$; compute $e f^{n}v$ on a highest-weight vector $v$ by
induction and find where the string $f^{n}v$ must terminate. P49: $\langle\varepsilon_i,\varepsilon_j\rangle = \delta_{ij}-\tfrac13$
is the trace form restricted to the traceless part; it gives
$|\varepsilon_i-\varepsilon_j|^2 = 2$. P51: cosine law: $\alpha_1\cdot\alpha_2 = (|\alpha_1|^2+|\alpha_2|^2-|\alpha_1-\alpha_2|^2)/2$. P52: for $\mathfrak{su}(3)$ the positive coroots are $\alpha_1,\alpha_2,\theta$ and
$\langle\rho,\cdot\rangle = 1,1,2$ on them. P53: the weights of $\mathbf3$ are $\omega_1$, $\omega_1-\alpha_1 = \omega_2-\omega_1$,
$\omega_1-\alpha_1-\alpha_2 = -\omega_2$; those of $\bar{\mathbf3}$ are
their negatives.

\section{Affine algebras and theta functions (Section~20)}

\begin{enumerate}[label=\textbf{P\arabic*.},leftmargin=*,start=55]
\item $(\star)$ With $t_\gamma$ acting at level $K$
  by $\mu\mapsto\mu+K\gamma$ on the finite part, show
  $w\,t_\gamma\,w^{-1} = t_{w\gamma}$, so that $W$ and $t(Q^\vee)$ generate
  $W\ltimes t(Q^\vee)$; show that $s_0 = t_{\theta^\vee}s_\theta$ is the
  reflection in the hyperplane $\langle\mu,\theta^\vee\rangle = K$; and for
  $\mathfrak{su}(2)$ at $K = k+2$ compute $s_0(m\omega) = (2(k+2)-m)\omega$.
  Conclude that the fundamental alcove is a simplex and that its interior
  lattice points at $K = k+h^\vee$ are exactly $P_k^+ + \rho$.
\item $(\star)$ Using
  $t_\gamma(\hat\mu) = \hat\mu+K\gamma-(\langle\mu,\gamma\rangle+\tfrac K2|\gamma|^2)\delta$
  and $e^{-\delta} = q$, show that
  $\sum_{\gamma\in Q^\vee}e^{t_\gamma(\hat\mu)} = q^{-|\mu|^2/2K}\,\Theta^{(K)}_\mu$
  with $\Theta^{(K)}_\mu$ as in Theorem~20.2, and conclude that
  the Weyl--Kac numerator is $\sum_{w\in W}\varepsilon(w)\Theta^{(K)}_{w(\lambda+\rho)}$
  up to a power of $q$ (why is the power the same for every $w$, and why
  is $\varepsilon(t_\gamma) = +1$?). State where in the computation the
  non-abelian structure enters and where the lattice does.
\item $(\star)$ From $h_\lambda = C_2(\lambda)/2(k+h^\vee)$,
  $c = k\dim\mathfrak g/(k+h^\vee)$ and the strange formula
  $|\rho|^2/2h^\vee = \dim\mathfrak g/24$, derive
  $T_\lambda = \exp2\pi i\big[|\lambda+\rho|^2/2(k+h^\vee)-\dim\mathfrak g/24\big]$.
  Specialize to $\mathfrak{su}(2)_k$; compute $c$ for $\mathfrak{su}(2)_1$,
  $\mathfrak{su}(2)_2$, $\mathfrak{su}(3)_1$ and $h_\lambda$ for the
  $\mathfrak{su}(2)_1$ doublet, and compare $T$ with $\rho_2(T)$ of the
  abelian register.
\item $(\star)$ (Frenkel--Kac.) Show that the two coset
  theta functions $\Theta_{Q+j\omega}$ of $A_1$ ($Q = \sqrt2\,\Z$) are the
  level-two functions $\theta^{(2)}_j$ of Definition~10.1 (with
  $z$ rescaled), and deduce $S$ and $T$ of $\mathfrak{su}(2)_1$ from the
  abelian register and the multiplier system of $\eta$
  ($\eta(\tau+1) = e^{\pi i/12}\eta(\tau)$, $\eta(-1/\tau) = \sqrt{-i\tau}\,\eta(\tau)$).
  Then do $\mathfrak{su}(3)_1$: the three cosets of the hexagonal lattice,
  $S$ the Fourier transform of $\Z/3$,
  $T\propto\operatorname{diag}(1,e^{2\pi i/3},e^{2\pi i/3})$. Explain why
  the even lattice $A_2$ has a diagonal $T$ while the rank-one register at
  odd level (P30) does not, and compare the
  $\mathfrak{su}(2)_1$ fusion rule with the Heisenberg--Weyl composition of
  characteristics.
\end{enumerate}

\paragraph{Hints.} P55: $\theta^\vee = \theta$ since
$\theta$ is long; the reflection in the affine hyperplane
$\langle\mu,\theta\rangle = K$ along $\theta$ is $\mu\mapsto\mu-(\langle\mu,\theta\rangle-K)\theta$. P56: $\tfrac K2|\tfrac\mu K+\gamma|^2 = \tfrac{|\mu|^2}{2K}+\langle\mu,\gamma\rangle+\tfrac K2|\gamma|^2$,
and $t_{\theta^\vee} = s_0s_\theta$. P57: for
$\mathfrak{su}(2)$, $\dim\mathfrak g = 3$, $|\rho|^2 = \tfrac12$,
$h^\vee = 2$; check the strange formula first. P58: for
$\gamma = \sqrt2(n+\tfrac j2)$, $|\gamma|^2/2 = (n+\tfrac j2)^2$, the
exponent of $q$ in $\theta^{(2)}_j$; for $A_2$,
$\langle\omega_1,\omega_1\rangle = \tfrac23$, $\langle\omega_1,\omega_2\rangle = \tfrac13$
and $\omega_2\equiv2\omega_1 \bmod Q$.

\section{Rank two, the alcove, fusion and density (Sections~21--25)}

\begin{enumerate}[label=\textbf{P\arabic*.},leftmargin=*,start=59]
\item $(\star)$ List the finite subgroups of $\GL(2,\Z)$ up to conjugacy and
  identify $D_2,D_3,D_4,D_6$ with $W(A_1\times A_1)$, $W(A_2)$, $W(B_2)$,
  $W(G_2)$ by exhibiting the reflections. (Compare P50.)
\item Determine the simplicial automorphism group of the $n\times n$ grid
  torus with the diagonal $(i,j)\!-\!(i{+}1,j{+}1)$, modulo translations, and of the triangular-
  lattice torus. Which is non-abelian? Explain why an automorphism can never
  realize a Dehn twist, and what the flip passes change so that the twists
  $T_A$, $T_B$ and the quarter turn $S$ are read on the square grid.
\item $(\star)$ Enumerate $P_k^+$ for $\mathfrak{su}(2)_k$ and
  $\mathfrak{su}(3)_k$ and count the states. Show the $\mathfrak{su}(2)$ count
  is $k+1$ and the $\mathfrak{su}(3)$ count is $(k+1)(k+2)/2$.
\item $(\star)$ Starting from Theorem~20.2, derive the explicit
  $\mathfrak{su}(2)$ character formula of Section~22, keeping
  track of the level shift $k\mapsto k+2$ and why the denominator has level
  $2h^\vee = 4$. Then check it numerically at $k = 1$: with $w = \tfrac12\langle\alpha,z\rangle$,
  the quotient $[\theta^{(6)}_{\lambda+1}(w)-\theta^{(6)}_{-(\lambda+1)}(w)]/[\theta^{(4)}_1(w)-\theta^{(4)}_{-1}(w)]$
  equals the Frenkel--Kac character $\theta^{(2)}_\lambda(w)/\eta(\tau)$
  (P58). Why must the check be made at generic $w$ rather than at $w = 0$?
\item Show that setting $W = 1$, $\rho = 0$, $h^\vee = 0$ in
  Theorem~20.2 recovers Definition~10.1: the abelian
  register is the no-root case. Why can this substitution not be made in
  the non-abelian ratio itself?
\item $(\star)$ For $\mathfrak{su}(2)_k$ verify that the Kac--Peterson $S$ is
  symmetric, unitary, and satisfies $S^2 = $ charge conjugation (here the
  identity). Compute $S$ and $T$ for $k = 1$ and $k = 2$. Show that with the
  $\dim\mathfrak g/24$ term in $T$ the relation $(ST)^3 = S^2$ holds exactly,
  not only projectively.
\item $(\star)$ Verify the $\mathfrak{su}(2)_k$ fusion rule of
  Theorem~23.2 from the Verlinde formula for $k = 2$, and
  reproduce it by Kac--Walton: add the weights of $V_\mu$ to $\lambda$ and
  reflect into the alcove. Compare with Proposition~19.2: which
  reflection is new?
\item $(\star)$ Prove Theorem~23.4 for $\mathfrak{su}(2)_k$: show
  $v_\lambda(\sigma) = S_{\lambda\sigma}/S_{0\sigma} = \sin\frac{(\lambda+1)(\sigma+1)\pi}{k+2}\big/\sin\frac{(\sigma+1)\pi}{k+2}$
  equals $\operatorname{ch}V_\lambda$ evaluated at $x = e^{-i\pi(\sigma+1)/(k+2)}$,
  and verify $v_1(\sigma)^2 = v_0(\sigma)+v_2(\sigma)$ pointwise at $k = 2$
  for $\sigma = 0,1,2$. Explain why the evaluation map is invertible and why
  this makes fusion ``multiplication of values''.
\item In the root coordinates of Section~22.1 with $A$ the (coroot) Gram
  matrix of $A_1$ or $A_2$: (a) show the characteristic set $P/KQ^\vee$ has
  $\det(A)K^r$ elements ($2K$ and $3K^2$); (b) show by Poisson summation that
  \[
  \Theta^{(K)}_\mu(-1/\tau,\xi/\tau) = (-i\tau)^{r/2}|P/KQ^\vee|^{-1/2}e^{\pi iK\xi^{\mathsf T}A\xi/\tau}
  \sum_{\mu'\in P/KQ^\vee}e^{-2\pi i\langle\mu,\mu'\rangle/K}\Theta^{(K)}_{\mu'}(\tau,\xi),
  \]
  and deduce that the exponential factors of numerator and denominator in the
  character ratio combine to $e^{\pi ik\xi^{\mathsf T}A\xi/\tau}$; (c) explain
  why $\Omega = \tau A$ is not a principal period matrix and why
  $(\tau A)\mapsto-(\tau A)^{-1}$ is not the transformation just derived.
\item $(\star)$ Explain, from the two theorems of Section~25,
  why the torus register is finite for every level and every algebra, and
  exactly what geometric ingredient (punctures, braids) universality requires.
  Map each ingredient onto a construction the code already has.
\item State precisely the hypothesis that turns the third claim of Section~1
  from algebra into geometry --- where the Weyl group acts, and what the
  surface and the algebra each supply --- and design the first experiment
  that could falsify it. (A torus with $D_6$ symmetry is a distractor: say why.)
\end{enumerate}

\paragraph{Hints.} P59: a finite subgroup preserves a positive definite
form, hence is a point group of the square or hexagonal lattice. P62: $\lambda+\rho = \lambda+1$ in units of the fundamental weight;
$|\alpha|^2 = 2$. P64: $\sum_{\mu}\sin\frac{\pi(\lambda+1)(\mu+1)}{k+2}
\sin\frac{\pi(\nu+1)(\mu+1)}{k+2}$ is a finite Fourier orthogonality. P66: $\operatorname{ch}V_\lambda(x) = (x^{\lambda+1}-x^{-\lambda-1})/(x-x^{-1})$ (P52). P67: Poisson summation in $n\in\Z^r$ for $f(x) = e^{\pi iK\tau x^{\mathsf T}Ax+2\pi iKx^{\mathsf T}A\xi}$; then group $m\in\Z^r$ into cosets modulo $KA\Z^2$.

\part*{The experiment specification}

\section{Whole harmonic states, pairings and elimination (Sections~26--28)}

\begin{enumerate}[label=\textbf{P\arabic*.},leftmargin=*,start=70]
\item $(\star)$ Prove the identifiability proposition of Section~26: with $A = R_{\rm in}Z$
  and $B = R_{\rm out}Z$, a unique operator $T$ with $B = TA$ defined on every
  input exists exactly when $A$ is onto and $\ker A\subseteq\ker B$, and then
  $T = BC$ for any right inverse $C$ of $A$, with the test $B(I-CA) = 0$.
  Give a $2\times2$ example where $A$ is not onto (an input that cannot be
  prescribed) and one where $\ker A\not\subseteq\ker B$ (an ambiguous output),
  and say which verifier fixtures (C2) exercise each.
\item For a scalar local system with edge factors $g_{ij}$ define
  $(d_Uf)(v_iv_j) = g_{ij}f(v_j)-f(v_i)$ and
  $(d_U\alpha)(v_0v_1v_2) = g_{01}\alpha(v_1v_2)-\alpha(v_0v_2)+\alpha(v_0v_1)$.
  Show $d_U^2f = (g_{01}g_{12}-g_{02})f(v_2)$ on a triangle, so $d_U^2 = 0$
  exactly when the holonomy around every face is trivial (flat), and that the
  natural pairing of $U$-cochains with chains needs the dual system $U^{-1}$
  for gauge invariance. Explain why the certificates of Section~26 read
  closedness on cochain images with $\partial^{U^{-1}}$ and coclosedness on the
  chain frame with $\partial^U$.
\item Derive the Gram conventions of Section~27. For boundary harmonic chain
  bases $\Phi,\Phi^\vee$ with period matrices $P,P^\vee$ of their images and
  $B_0 = (\Phi^\vee)^{\mathsf T}M_{\rm own}\Phi$, show that the bilinear pairing
  in period coordinates is $G_{\rm per} = (P^\vee)^{-\mathsf T}B_0P^{-1}$, that the
  abbreviation $(P^\vee)^{-\mathsf T}P^{-1}$ requires $B_0 = I$, and that with
  primal/dual transports $M,M^\vee$ pairing preservation is
  $(M^\vee)^{\mathsf T}G_BM = G_A$. Using Remark~7.1, show the pairing is gauge
  invariant while a primal--primal Gram is not.
\item $(\star)$ Prove the elimination identities of Section~28: for
  $K = \begin{psmallmatrix}K_{bb}&K_{bz}\\K_{zb}&K_{zz}\end{psmallmatrix}$ with
  $K_{zz}$ invertible, $Kx = 0$ forces $z = -K_{zz}^{-1}K_{zb}b$ and
  $S_Gb = 0$ with $S_G = K_{bb}-K_{bz}K_{zz}^{-1}K_{zb}$. Then minimize
  $\|u\|^2+\|r-T_2u\|^2$ over $u$ and show the minimum is
  $r^\dagger(I+T_2T_2^\dagger)^{-1}r$, attained at $u = (I+T_2^\dagger T_2)^{-1}T_2^\dagger r$.
  Why is this exact elimination and not a partial trace?
\end{enumerate}

\paragraph{Hints.} P70: $c-CAc\in\ker A$ for every $c$. P73: Woodbury: $I-T_2(I+T_2^\dagger T_2)^{-1}T_2^\dagger = (I+T_2T_2^\dagger)^{-1}$.

\section{Quantum norms, channels and acceptance (Sections~29--30)}

\begin{enumerate}[label=\textbf{P\arabic*.},leftmargin=*,start=74]
\item $(\star)$ With declared positive $Q_{\rm in},Q_{\rm out}$ and
  $\widehat T = Q_{\rm out}^{1/2}TQ_{\rm in}^{-1/2}$, show
  $T^\dagger Q_{\rm out}T = Q_{\rm in}$ iff $\widehat T^\dagger\widehat T = I$;
  that the reference marginal of $|\widehat T\rangle\!\rangle$ is
  $(\widehat T^\dagger\widehat T)^{\mathsf T}/\|\widehat T\|_F^2$, equal to
  $I/d_{\rm in}$ for an isometry; and that the Schmidt coefficients of the
  normalized operator vector are the normalized singular values of $\widehat T$.
\item For Kraus operators $K_a$ show that $J = \sum_a|K_a\rangle\!\rangle\langle\!\langle K_a|$
  is positive and that $\Tr_{\rm out}J = I_{\rm in}$ iff $\sum_aK_a^\dagger K_a = I$;
  build the Stinespring isometry $V = \sum_aK_a\otimes|a\rangle_E$ and check
  $V^\dagger V = I$, $\mathcal E(\rho) = \Tr_E(V\rho V^\dagger)$. Then explain by a
  dimension count why a third boundary component, which enters the relation
  as a direct summand $B\oplus E$, cannot play the role of the tensor factor
  $B\otimes E$ of a dilation.
\item Section~29 distinguishes three tensor constructions. (a) Compute
  $H^2(T^2\times T^2)$ by K\"unneth and identify the summand occupied by the
  degree-one registers of the two factors. (b) Give the dimension of the
  theta register for a genus-$g$ surface at level $k$ and for two disjoint
  tori. (c) For $\mathfrak{su}(2)_1$ compare $V_1\otimes V_1$ with the fusion
  $1\times1$. (d) Two level-two theta factors give four amplitudes; count the
  amplitudes of a two-qubit gate, of its operator vector, of a direct sum of
  two two-dimensional port frames, and of a $2\times2$ selected-pair $\chi$,
  and explain why matching counts do not define subsystems.
\item V1--V8 require each check to declare whether it compares a geometric
  computation with an independent oracle or tests a certificate. Classify
  the verifier's families R1--R5, H1--H3, J1--J3, C1, D1, U1, T1--T9, F1--F4,
  G1--G5 and E1--E6 accordingly, naming the oracle or the certified statement.
\end{enumerate}

\paragraph{Hints.} P75: use P74 for the marginal of each $|K_a\rangle\!\rangle$.

\section{Neck geometry and encoded-state entropy (Section~33)}

\begin{enumerate}[label=\textbf{P\arabic*.},leftmargin=*,start=78]
\item $(\star)$ Prove Proposition~33.2. (a) If the smaller Schmidt probability of a
  two-qubit state is $p = C\epsilon^2+O(\epsilon^3)$, show
  $S = -C\epsilon^2\log_2(C\epsilon^2)+C\epsilon^2/\ln2+o(\epsilon^2)$, so the
  entropy is not proportional to $\epsilon^2$. (b) For
  $\Omega = \operatorname{diag}(\tau_1,\tau_2)+\epsilon\begin{psmallmatrix}0&1\\1&0\end{psmallmatrix}$
  show $\partial_\epsilon\theta_j|_{\epsilon=0} = \theta'_{j_1}(z_1;\tau_1)\theta'_{j_2}(z_2;\tau_2)/(2\pi ik)$,
  that at $z = 0$ this vanishes for every level-two characteristic, and
  hence that $p\sim C\epsilon^4$ there. (c) Describe a numerical check of (b)
  with \texttt{ThetaRegister.coherent\_state} and \texttt{entanglement\_entropy}.
\end{enumerate}

\paragraph{Hints.} P78: level-two $\theta_j$ are even in $z$: substitute $n\mapsto-n$ ($j = 0$) or $n\mapsto-n-1$ ($j = 1$).

\part*{The measured constructions}

\section{Reading the records (Section~35)}

\begin{enumerate}[label=\textbf{P\arabic*.},leftmargin=*,start=79]
\item $(\star)$ Run the verifier on the collar fixtures,
\begin{verbatim}
python examples/cobordism/qubit_animation.py verify --tori 2
python examples/cobordism/qubit_animation.py verify --tori 4 \
    --collar-twist swap
python examples/cobordism/qubit_animation.py theta --tori 4
\end{verbatim}
  read the records (Section~35.2 names their directories), and for every check name
  the statement of the reading it measures. Which checks would still pass if
  the harmonic band were replaced by an arbitrary two-dimensional subspace
  of closed cochains, and which would fail?
\item $(\star)$ The flip passes are declared by their flipped edge, tetrahedron, new edge
  and lattice map (\texttt{DECLARED\_FLIP\_PASSES}). Pass \texttt{b} flips the
  row edge $(1,0)$ inside the quadrilateral with vertices $(0,-1)$, $(0,0)$,
  $(1,0)$, $(1,1)$ to the new edge from $(0,-1)$ to $(1,1)$; the three lattice
  maps are
  \[
   L_b = \begin{pmatrix}1&0\\1&1\end{pmatrix},\qquad
   L_a = \begin{pmatrix}1&1\\0&1\end{pmatrix},\qquad
   L_s = \begin{pmatrix}0&1\\-1&0\end{pmatrix}.
  \]
  (a) Show that the edge set of the flipped grid is the image under $L_b$ of
  the grid's edge set, the displacements $(1,0)$, $(0,1)$ and $(1,1)$, and that
  the far marking read through $L_b$ gives $M = L_b^{\mathsf T} = T_B$; do the
  same for $L_a$ and $L_s$.
  (b) Show that two \texttt{b} passes on the $3\times3$ grid produce a new edge
  whose vertex pair is already joined by an edge of the host, while the
  $5\times5$ grid separates them. (c) Run
\begin{verbatim}
python examples/cobordism/qubit_animation.py verify --far-flips b
python examples/cobordism/qubit_animation.py verify --grid 5 --far-flips ab
python examples/cobordism/qubit_animation.py verify --grid 5 --far-flips ba
\end{verbatim}
  and confirm F2--F4 and the modulus relation F3 of P36.
\item The tube host joins two collars by a prism tube between one far triangle
  of each. (a) Show that the far boundary is a genus-two surface and that
  $\chi(\partial W) = -2$, then use $\chi(W) = \tfrac12\chi(\partial W)$ and
  Mayer--Vietoris to derive the Betti numbers $[1,4,2,0]$ and the restriction
  rank $4 = \tfrac12b_1(\partial W)$. (b) Explain why the monodromy from the two
  near tori to $\Sigma_2$ is block-diagonal (G2) although the period matrix of
  $\Sigma_2$ has $\Omega_{12}\ne0$ (G3). (c) With
\begin{verbatim}
python examples/cobordism/qubit_animation.py theta --tori 4 --join tube
python examples/cobordism/qubit_animation.py theta --tori 4 --join tube \
    --neck-waists 0.25,0.5,1,2,4 --neck-lengths 0.25,0.5,1,2,4
\end{verbatim}
  record $\Omega$, the coherent-state entropy and the trend of $|\Omega_{12}|$
  with the neck, and explain the trend by the plumbing description of a
  degenerating surface.
\item $(\star)$ Build the $\Z/10$-symmetric genus-two surface,
\begin{verbatim}
from tessera.quantum import SymmetricGenusTwo, ThetaRegister
import numpy as np
s = SymmetricGenusTwo()
\end{verbatim}
  and check: $28$ vertices, $90$ edges, $60$ faces, $\chi = -2$; \texttt{s.rotation(1)}
  is a simplicial automorphism (\texttt{s.is\_automorphism}); $M = $ \texttt{s.induced\_map(2)}
  is the matrix of P47 and satisfies $M^{\mathsf T}JM = J$, $M^5 = I$. Then
  fit $\rho_2(M)$ with \texttt{ThetaRegister(level=2).weil(M, s.periods.omega)},
  verify unitarity and $\rho^5\propto I$, compute the operator Schmidt rank
  across the two factors, and count how many of the $36$ product stabilizer
  states are sent to states of exactly one bit of entanglement. Compare with
  \texttt{theta --surface decagon --rotation 2}.
\item \texttt{SurfacePeriods(faces, lengths, a\_cycles, b\_cycles)} reads a period
  matrix from a triangulated surface: cotangent weights, harmonic $1$-cochains
  as the kernel of $[d_1;\,d_0^{\mathsf T}\operatorname{diag}(w)]$, Whitney
  interpolation, the Hodge star $J$ with $J^2 = -I$, the $-i$ eigenspace and
  $\Omega = P_BP_A^{-1}$. (a) On a flat square-grid torus of modulus $\tau$
  check that it returns $\Omega = \tau$ and that the cup pairing gives
  $A\cdot B = +1$. (b) For the decagon surface read $\Omega$ from
  \texttt{s.periods}, check symmetry, positivity of $\operatorname{Im}\Omega$
  (eigenvalues $0.46$, $1.66$), and that $\Omega$ is fixed by $M$ under the
  bridge of P36: $(R+S\Omega)(P+Q\Omega)^{-1} = \Omega$. Why must a
  surface with an order-five symplectic automorphism have a fixed period
  matrix, and why does this not make $\Omega = \tau A$ for a root Gram $A$?
\end{enumerate}

\paragraph{Hints.} P80: the new edge of the second pass has old-coordinate displacement $L_b(1,2)^{\mathsf T} = (1,3)^{\mathsf T}$. P81: a $1$-handle attached to two tori along disks: $\chi = 0+0-2$. P82: the product stabilizer states are $|a\rangle\otimes|b\rangle$ with $a,b$ among the six eigenstates of $X,Y,Z$. P83: an automorphism acts on $H^1$ by an integer symplectic matrix and on the holomorphic forms by a change of basis.

\end{document}
